\documentclass[10pt,a4paper]{ctexart}

\usepackage[left=0.5in,right=0.5in,top=0.5in,bottom=0.5in]{geometry}
\usepackage{hyperref}
\usepackage{enumitem}
\usepackage{parskip}
\usepackage{xcolor}
\pagenumbering{gobble}

\definecolor{darkred}{RGB}{153,0,0}

\hypersetup{
    colorlinks=true,
    linkcolor=darkred,
    urlcolor=darkred,
    pdftitle={王子翔-实习简历},
    pdfauthor={王子翔}
}

\newcommand{\sectiontitle}[1]{
    \vspace{4mm}
    {\large\textbf{#1}}
    \vspace{1mm}
    \hrule
    \vspace{3mm}
}

\newcommand{\entryhead}[2]{
    {#1} \hfill #2 \\
}

\newcommand{\projectentry}[4]{
    \noindent
    \textbf{#1} \hfill #2 \\
    {\small #3} \\
    #4
    \vspace{2mm}
}

\newcommand{\pubentry}[4]{
    \noindent
    \textbf{#1} \hfill #2 \\
    #3 \\
    {\small #4}
    \vspace{2mm}
}

\setlength{\parindent}{0pt}
\setlength{\parskip}{0.5mm}
\setlist[itemize]{leftmargin=*,topsep=2pt,itemsep=1pt,parsep=0pt}

\begin{document}

\begin{center}
{\huge\textbf{王子翔}}\\
\vspace{2mm}
\textbf{电话：} 18949871109 \hspace{4mm}
\textbf{邮箱：} \href{mailto:wangzx@stu.pku.edu.cn}{wangzx@stu.pku.edu.cn} \hspace{4mm}
\textbf{主页：} \href{https://fieldry.github.io}{fieldry.github.io}
\end{center}

\sectiontitle{求职意向}

算法实习 / 研究实习（机器学习、大模型、数据挖掘、AI Agent、医疗 AI 方向可灵活适配）

\sectiontitle{教育背景}

\entryhead{\textbf{北京大学}，软件与微电子学院，软件工程硕士}{2024.09 - 至今}
\begin{itemize}
    \item 研究方向：AI for Healthcare、医疗大语言模型、电子病历数据建模
    \item 导师：马连韬
    \item 可根据投递岗位补充：核心课程、GPA、排名、研究课题关键词
\end{itemize}

\entryhead{\textbf{北京航空航天大学}，软件学院，软件工程学士}{2019.09 - 2024.06}
\begin{itemize}
    \item 系统掌握软件工程、数据结构、算法设计、系统实现等基础能力
    \item 可根据投递岗位补充：课程成绩、竞赛经历、工程项目
\end{itemize}

\sectiontitle{个人概述}

我目前主要从事\textbf{大语言模型与数据建模相关研究}，关注\textbf{多智能体协作、检索增强、结构化时序数据建模、模型评测与工作流自动化}。在项目中，我更偏向承担\textbf{问题建模、方法设计、实验实现、代码复现与结果分析}等工作。后续可按岗位需求进一步强化为“算法研究”“LLM 应用”“数据挖掘”或“工程实现”版本。

\sectiontitle{研究 / 项目经历}

\projectentry
{ColaCare：基于多智能体协作的电子病历建模}
{2024 - 2025}
{关键词：LLM Agent、多智能体协作、医疗数据建模、实验设计}
{
\begin{itemize}
    \item 参与面向电子病历场景的多智能体协作框架研究，围绕任务拆解、交互流程与性能提升开展实验。
    \item 负责或参与模型实验、结果分析、论文整理等工作；可按后续实际贡献补充更具体的技术细节。
    \item 对外成果发表于 WWW 2025，可作为研究能力与论文产出代表项目。
\end{itemize}
}

\projectentry
{RetCare / EHRFlow：医疗任务中的 LLM 检索与工作流自动化}
{2024}
{关键词：RAG、Medical LLM、Workflow、Prompting}
{
\begin{itemize}
    \item 参与围绕医疗知识检索、任务流程设计和自动化分析的项目研究。
    \item 涉及模型调用、实验配置、评测流程和结果整理等工作。
    \item 可按投递方向将表述切换为“Agent 工作流设计”“Prompt 与检索增强”“实验平台搭建”等版本。
\end{itemize}
}

\projectentry
{EHR 预测建模与评测基准}
{2023 - 至今}
{关键词：时间序列 EHR、临床预测、缺失值处理、Benchmark}
{
\begin{itemize}
    \item 持续参与电子病历时序数据建模、预测任务评测与协议整理等工作。
    \item 相关成果覆盖 PRISM、EMERGE、Protocol、Benchmark 等项目，兼具方法研究与规范化数据处理经验。
    \item 可针对岗位重点突出“表格时序建模”“数据清洗与特征工程”“训练评测 pipeline”能力。
\end{itemize}
}

\sectiontitle{实习经历（待补充）}

\projectentry
{实习经历占位 1}
{时间待补充}
{公司 / 团队 / 岗位待补充}
{
\begin{itemize}
    \item 建议补充 3 条以内高质量 bullet，优先描述目标、方法、结果和指标。
    \item 尽量量化：模型效果提升、系统吞吐、数据规模、业务指标、上线范围等。
    \item 若涉及大模型应用，建议明确 prompt、agent、RAG、评测、部署中的具体职责。
\end{itemize}
}

\projectentry
{实习经历占位 2}
{时间待补充}
{公司 / 团队 / 岗位待补充}
{
\begin{itemize}
    \item 若后续有工程向经历，可强调数据 pipeline、训练框架、推理服务、评测平台等内容。
    \item 若后续有研究向经历，可强调课题问题定义、方法创新、实验设计和论文产出。
\end{itemize}
}

\sectiontitle{代表成果}

\pubentry
{ColaCare: Enhancing Electronic Health Record Modeling through Large Language Model-Driven Multi-Agent Collaboration}
{WWW 2025}
{\textbf{王子翔}$^{\ast}$，朱英豪$^{\ast}$，赵慧雅$^{\ast}$ 等}
{多智能体协作与 EHR 建模代表成果}

\pubentry
{ClinicRealm: Re-evaluating Large Language Models with Conventional Machine Learning for Non-Generative Clinical Prediction Tasks}
{npj Digital Medicine 2026}
{朱英豪，高峻逸，\textbf{王子翔} 等}
{围绕 LLM 与传统方法在临床预测任务中的比较与再评估}

\pubentry
{Prompting Large Language Models for Zero-Shot Clinical Prediction with Structured Longitudinal Electronic Health Record Data}
{Preprint 2024}
{朱英豪$^{\ast}$，\textbf{王子翔}$^{\ast}$，高峻逸 等}
{结构化时序 EHR 与零样本临床预测相关工作}

\sectiontitle{技术栈}

\begin{itemize}
    \item \textbf{编程与框架：} Python、PyTorch、基础机器学习 / 深度学习实验开发
    \item \textbf{研究与算法：} 大语言模型应用、Prompting、RAG、多智能体协作、时序 / 表格数据建模
    \item \textbf{工程能力：} 数据处理、实验管理、结果分析、论文与技术文档整理
    \item \textbf{后续可补充：} C++ / Java / SQL / Linux / Docker / 分布式训练 / 服务部署等与岗位匹配的技能项
\end{itemize}

\sectiontitle{荣誉与服务}

\begin{itemize}
    \item 北京航空航天大学优秀毕业生（2024）
    \item 国家励志奖学金（2022，2023）
    \item 课程助教：数据结构与程序设计、离散数学、计算机硬件基础
    \item 会议审稿：AAAI 2026
\end{itemize}

\sectiontitle{补充说明}

这份简历当前是\textbf{面向找实习的基础框架}，后续建议优先补充：
\begin{itemize}
    \item 真实实习经历及量化结果
    \item 具体工程项目与代码链接
    \item 可公开展示的技术栈、比赛、开源贡献
    \item 针对不同岗位版本的关键词适配
\end{itemize}

\end{document}
